% from ega4-17-fr.tex
\oldpage[IV-4]{69}
L'équivalence de la condition c) de (17.5.1) et de d) résulte de l'équivalence de a) et d) dans le critère jacobien (0, 22.6.4), et l'équivalence de d) et d') résulte de (0, 19.3.6). D'autre part, d'') implique d') sans hypothèse sur les corps résiduels (0, 19.3.4). Enfin, si $\mathfrak m$ désigne l'idéal maximal de $\widehat A$, l'hypothèse d') entraîne que $\widehat B/\mathfrak m\widehat B$ est une $k(y)$-algèbre locale noethérienne complète, formellement lisse pour sa topologie adique (0, 19.3.5); l'hypothèse $k(y)=k(x)$ entraîne donc que $\widehat B/\mathfrak m\widehat B$ est $k(y)$-isomorphe à une algèbre de séries formelles $k(y)[[T_1,\ldots,T_n]]$ (0, 19.6.4). Comme d'autre part, $\widehat A[[T_1,\ldots,T_n]]$ est un $\widehat A$-module plat et une $\widehat A$-algèbre locale noethérienne complète, on conclut de (0, 19.7.1.5) que cette algèbre est isomorphe à $\widehat B$. Donc d') entraîne d'') sous l'hypothèse additionnelle $k(x)=k(y)$.

\begin{remark}[17.5.4]
\label{IV.17.5.4}
Supposons que $Y$ soit un préschéma localement noethérien, et $f:X\to Y$ un morphisme localement de type fini. Le critère (17.5.3, d)), joint à (0, 22.1.4) montre que pour démontrer que $f$ est lisse, on peut appliquer la définition (17.1.1), en se restreignant au cas où le schéma affine $Y'$ est le spectre d'un anneau local artinien.
\end{remark}

\begin{proposition}[17.5.5]
\label{IV.17.5.5}
Soient $Y$ un préschéma localement noethérien, $f:X\to Y$ un morphisme localement de type fini, $x$ un point de $X$, $y=f(x)$. Supposons que $Y$ soit réduit au point $y$. Alors, pour que $f$ soit lisse au point $x$, il faut et il suffit que $f$ soit universellement ouvert dans un voisinage de $x$ dans $f^{-1}(y)$ et que $f^{-1}(y)$ soit un $k(y)$-préschéma géométriquement régulier au point $x$.
\end{proposition}

Compte tenu de (17.5.1), tout revient à voir que, si $f^{-1}(y)$ est un $k(y)$-préschéma géométriquement régulier au point $x$, il est équivalent de dire que $f$ est plat au point $x$ ou universellement ouvert dans un voisinage de $x$ dans $f^{-1}(y)$. Or, si $f$ est plat au point $x$, il l'est dans un voisinage de $x$ dans $X$ (11.1.1) et par suite est universellement ouvert dans ce voisinage (2.4.6). Réciproquement, l'hypothèse que $f$ est universellement ouvert dans un voisinage de $x$ dans $f^{-1}(y)$ et que $f^{-1}(y)$ est un $k(y)$-préschéma géométriquement régulier au point $x$ entraîne que $f$ est plat au point $x$, puisque $\sh{O}_{Y,y}$ est réduit (15.2.2).

\begin{corollary}[17.5.6]
\label{IV.17.5.6}
Soient $Y$ un préschéma localement noethérien, $f:X\to Y$ un morphisme localement de type fini, $x$ un point de $X$, $y=f(x)$. On suppose que $Y$ soit réduit et géométriquement unibranche (6.15.1) au point $y$. Alors, pour que $f$ soit lisse au point $x$, il faut et il suffit que $f$ soit équidimensionnel au point $x$ et que $f^{-1}(y)$ soit un $k(y)$-préschéma géométriquement régulier au point $x$.
\end{corollary}

En remarquant que l'ensemble des points où $f$ est équidimensionnel est ouvert (13.3.2), on voit que le corollaire résulte de (17.5.5) et du critère de Chevalley (14.4.4).

Le fait que $f$ soit un morphisme lisse en un point entraîne en particulier que $f$ vérifie

